\documentclass[11pt]{article}

\usepackage[margin=1in]{geometry}
\usepackage{amsmath, amssymb}
\usepackage{enumitem}
\usepackage{array}
\usepackage{fancyhdr}

\pagestyle{fancy}
\fancyhf{}
\lhead{Scientific Notation Practice (10.5--10.7)}
\rhead{Name:\ \rule{2.7in}{0.4pt}}
\cfoot{\thepage}

\setlength{\parindent}{0pt}

\newcommand{\work}[1]{\vspace{#1}}

\begin{document}

\begin{center}
{\Large \textbf{Practice Worksheet \#2: Sections 10.5--10.7}}\\[4pt]
{\small Reading, Writing, and Operating with Scientific Notation}
\end{center}

\textbf{Directions:} Show all work. For Section D, write each final answer in \textbf{scientific notation}.

\vspace{0.4em}
\hrule
\vspace{0.9em}

% =========================================================
% 10.5 Reading Scientific Notation
% =========================================================
\textbf{A. Tell whether the number is written in scientific notation. Explain. (Section 10.5)}

\begin{enumerate}
\item $4.2 \times 10^{9}$
\work{1in}

\item $0.58 \times 10^{3}$
\work{1in}

\item $9.87 \times 10^{-2}$
\work{1in}

\item $15.6 \times 10^{4}$
\work{1in}
\end{enumerate}

\vspace{0.2em}
\hrule
\vspace{0.9em}

% =========================================================
% 10.5 Standard Form
% =========================================================
\textbf{B. Write the number in standard form. (Section 10.5)}

\begin{enumerate}[resume]
\item $7.25 \times 10^{6}$
\work{1.1in}

\item $4.01 \times 10^{-3}$
\work{1.1in}

\item $2.9 \times 10^{1}$
\work{1.1in}

\item $6.84 \times 10^{-5}$
\work{1.1in}
\end{enumerate}

\vspace{0.2em}
\hrule
\vspace{0.9em}

% =========================================================
% 10.6 Writing Scientific Notation
% =========================================================
\textbf{C. Write the number in scientific notation. (Section 10.6)}

\begin{enumerate}[resume]
\item $0.00000672$
\work{1.1in}

\item $81,\!400,\!000$
\work{1.1in}

\item $0.0935$
\work{1.1in}

\item $6,\!020,\!000$
\work{1.1in}
\end{enumerate}

\vspace{0.2em}
\hrule
\vspace{0.9em}

% =========================================================
% 10.7 Operations
% =========================================================
\textbf{D. Evaluate the expression. Write your answer in scientific notation. (Section 10.7)}

\begin{enumerate}[resume]
\item $(4.8\times10^{4}) + (3.6\times10^{4})$
\work{1.5in}

\item $(6.2\times10^{-5}) - (1.9\times10^{-6})$
\work{1.5in}

\item $(7.5\times10^{2}) \times (4.0\times10^{-3})$
\work{1.7in}

\item $(9.6\times10^{-6}) \div (1.2\times10^{-9})$
\work{1.7in}

\item $(3.4\times10^{8}) + (6.5\times10^{7})$
\work{1.7in}

\item $(8.1\times10^{-2}) - (2.7\times10^{-3})$
\work{1.7in}
\end{enumerate}

\vspace{0.2em}
\hrule
\vspace{0.9em}

% =========================================================
% Applications
% =========================================================
\textbf{E. Applications (Sections 10.5--10.7)}

\textbf{1. PLANETS.} The table shows approximate equatorial radii of several planets.

\vspace{0.4em}
\begin{tabular}{|>{\raggedright\arraybackslash}p{2.1in}|>{\raggedleft\arraybackslash}p{2.6in}|}
\hline
\textbf{Planet} & \textbf{Equatorial Radius (km)}\\
\hline
Mercury & $2.44\times10^{3}$\\
Venus   & $6.05\times10^{3}$\\
Earth   & $6.38\times10^{3}$\\
Mars    & $3.40\times10^{3}$\\
Jupiter & $7.15\times10^{4}$\\
Saturn  & $6.03\times10^{4}$\\
\hline
\end{tabular}

\begin{enumerate}[label=\alph*.]
\item Which planet has the \textbf{smallest} equatorial radius? Explain.
\work{1.3in}

\item Which planet has the \textbf{largest} equatorial radius? Explain.
\work{1.3in}
\end{enumerate}

\vspace{0.6em}

\textbf{2. DISTANCE.} A satellite is $7\times10^{6}$ kilometers from Earth.  
If $1$ kilometer $=10^{3}$ meters, how far is the satellite from Earth in meters?
Write your answer in scientific notation.
\work{1.8in}

\textbf{3. BIOLOGY.} A cell membrane is $0.00000091$ meters thick.
Write this number in scientific notation.
\work{1.4in}

\textbf{4. ORBITS.} The Sun takes about $2.4\times10^{8}$ years to orbit the Milky Way.
A planet takes $1.5\times10^{1}$ years to orbit its star.
How many times does the planet orbit while the Sun completes one orbit?
Write your answer in standard form.
\work{2.0in}

\end{document}